\begin{table}[H]
\centering
\caption{Crecimiento de la remuneración laboral por departamento, 2010--2025}
\label{tab:dept_remuneracion_crecimientos_24}
\scriptsize
\begin{tabular}{lrrrr}
\toprule
Departamento & Crec. rem./trab. & Puesto & Crec. rem./hora & Puesto \\
\midrule
Cundinamarca & 2,9\% & 1 & 3,2\% & 1 \\
Risaralda & 2,7\% & 2 & 3,1\% & 2 \\
Caldas & 1,9\% & 3 & 2,5\% & 3 \\
Boyacá & 1,9\% & 4 & 2,2\% & 6 \\
Tolima & 1,8\% & 5 & 2,1\% & 8 \\
Valle del Cauca & 1,8\% & 6 & 2,2\% & 5 \\
Antioquia & 1,7\% & 7 & 2,1\% & 9 \\
Cauca & 1,7\% & 8 & 2,0\% & 10 \\
Bogotá D.C. & 1,6\% & 9 & 2,2\% & 7 \\
Quindío & 1,4\% & 10 & 2,2\% & 4 \\
Chocó & 1,3\% & 11 & 1,9\% & 12 \\
Caquetá & 1,3\% & 12 & 1,7\% & 17 \\
Bolívar & 1,3\% & 13 & 1,9\% & 11 \\
Meta & 1,2\% & 14 & 1,7\% & 16 \\
Atlántico & 1,2\% & 15 & 1,7\% & 15 \\
Huila & 1,1\% & 16 & 1,1\% & 23 \\
Santander & 1,1\% & 17 & 1,3\% & 22 \\
Norte de Santander & 1,0\% & 18 & 1,4\% & 21 \\
Córdoba & 1,0\% & 19 & 1,7\% & 14 \\
Nariño & 0,9\% & 20 & 1,8\% & 13 \\
Cesar & 0,7\% & 21 & 1,4\% & 19 \\
Magdalena & 0,6\% & 22 & 1,5\% & 18 \\
Sucre & 0,3\% & 23 & 1,4\% & 20 \\
La Guajira & -0,9\% & 24 & -0,7\% & 24 \\
\bottomrule
\end{tabular}
\caption*{\footnotesize Nota: tasas de crecimiento anualizadas. Los indicadores de remuneración se calculan entre ocupados con ingreso horario positivo y horas válidas. Fuente: cálculos propios con GEIH.}
\end{table}